% gpu-nvme-direct: GPU-Initiated NVMe I/O on Consumer Hardware
% Draft v0.2 — 2026-02-28

\documentclass{article}
\usepackage{arxiv}
\usepackage[utf8]{inputenc}
\usepackage[T1]{fontenc}
\usepackage{microtype}
\usepackage{graphicx}
\usepackage{booktabs}
\usepackage{tabularx}
\usepackage{listings}
\usepackage{xcolor}
\usepackage{hyperref}
\usepackage{url}
\usepackage{amsmath}
\usepackage{subcaption}
\usepackage{enumitem}
\hypersetup{
  colorlinks=true,
  linkcolor=blue,
  citecolor=blue,
  urlcolor=blue
}

% Code listing style
\lstset{
  language=C,
  basicstyle=\ttfamily\footnotesize,
  keywordstyle=\color{blue}\bfseries,
  commentstyle=\color{gray},
  stringstyle=\color{red},
  breaklines=true,
  frame=single,
  numbers=left,
  numberstyle=\tiny\color{gray},
  captionpos=b,
  aboveskip=8pt,
  belowskip=8pt,
}

\newcommand{\sysname}{\textsc{gpu-nvme-direct}}
\newcommand{\us}{\,\text{\textmu s}}
\newcommand{\MBs}{\,\text{MB/s}}
\newcommand{\GBs}{\,\text{GB/s}}

\title{\sysname{}: GPU-Initiated NVMe I/O on Consumer Hardware}

\author{
  Samuel Cortes \\
  \texttt{https://naranjositos.tech/} \\
  \texttt{https://github.com/xaskasdf/gpu-nvme-direct}
}

\date{February 2026}

\begin{document}

\maketitle

% ============================================================
% ABSTRACT
% ============================================================
\begin{abstract}

We present \sysname{}, a system that enables a consumer GPU (NVIDIA GeForce RTX
3090) to autonomously initiate NVMe storage I/O operations via PCIe MMIO,
entirely without CPU involvement in the data path. Unlike prior work that
requires enterprise GPUs with native peer-to-peer (P2P) DMA support,
\sysname{} operates on commodity hardware (AMD Ryzen 5800X, B450 motherboard)
costing approximately \$2{,}000, using a \emph{tiered approach} that degrades
gracefully based on available P2P capabilities.

The system works by having a single CUDA thread construct NVMe submission queue
entries in host pinned memory, write doorbell registers on the NVMe
controller's BAR0 via PTX \texttt{st.relaxed.mmio.sys} instructions, and poll
completion queue entries---all without CPU intervention. This is possible
because PCIe \emph{posted} Memory Write TLPs succeed through the AMD data
fabric even on consumer platforms where non-posted reads fail.

We evaluate on two NVMe devices: a WD SN530 (PCIe~3.0~$\times$4), where
\sysname{} achieves 2{,}666\MBs{} (78\% of link bandwidth), and a WD SN740
(PCIe~4.0~$\times$4), where it reaches 3{,}350\MBs{} sustained. Compared to
CPU-mediated baselines, \sysname{} delivers up to 2.2$\times$ higher
throughput due to its ability to maintain 32 in-flight NVMe commands from a
single GPU thread. We demonstrate practical applicability for LLM inference:
integrated into our streaming inference engine, \sysname{} achieves
0.06~tok/s for Llama~70B (5$\times$ over mmap+memcpy), and combined with
tiered caching reaches 0.5~tok/s on a single RTX~3090.

\end{abstract}

\keywords{GPU-initiated I/O \and NVMe \and PCIe peer-to-peer \and consumer hardware \and LLM inference \and storage access}

% ============================================================
% 1. INTRODUCTION
% ============================================================
\section{Introduction}
\label{sec:intro}

The explosive growth of large language models (LLMs) has created a fundamental
tension between model size and available GPU memory. Models such as Llama~2
70B~\cite{touvron2023llama2} require $\sim$70\,GB in FP16, or $\sim$42\,GB in
Q6\_K quantization---far exceeding the 24\,GB VRAM of consumer GPUs like the
NVIDIA RTX 3090. Running these models on consumer hardware requires
\emph{streaming} model weights from storage to GPU memory, one layer at a time.

The conventional data path for this streaming is CPU-mediated:
\begin{center}
\small
NVMe $\xrightarrow{\text{DMA}}$ Page Cache $\xrightarrow{\text{CPU memcpy}}$
Pinned Memory $\xrightarrow{\text{H2D DMA}}$ GPU VRAM
\end{center}

This pipeline introduces a fundamental bottleneck: the CPU must copy every byte
of model data through its caches, consuming memory bandwidth and CPU cycles
that could be used for other tasks. In practice, this limits throughput to
$\sim$1.5--2\GBs{} and yields only 0.02~tokens/s for 70B models on consumer
hardware.

NVIDIA's GPUDirect Storage (GDS)~\cite{gds} addresses this by enabling direct
DMA from NVMe to GPU memory, bypassing the CPU. However, GDS requires
enterprise GPUs (Tesla, A-series, H-series) and certified storage hardware---a
software stack explicitly unavailable on GeForce consumer GPUs. Prior academic
work such as BaM~\cite{qureshi2023bam} demonstrated that GPUs can
\emph{initiate} NVMe I/O operations autonomously, but again requires enterprise
hardware with native PCIe peer-to-peer support.

In this paper, we ask: \emph{Can a consumer GPU act as an autonomous I/O
processor, directly driving an NVMe controller without any CPU involvement in
the data path?}

We answer this question affirmatively with \sysname{}, a system that enables
GPU-initiated NVMe I/O on commodity hardware. Our key insight is that even on
platforms where PCIe P2P reads fail (AMD consumer root complexes silently drop
non-posted read TLPs), \emph{posted} PCIe Memory Write TLPs succeed---and
writes are all that is needed for a GPU to drive an NVMe controller. The GPU
writes NVMe command entries to host pinned memory, writes doorbell registers on
the NVMe BAR0 via MMIO, and polls completion entries from host pinned memory.
The NVMe controller independently DMAs data to/from host memory. The CPU is
entirely uninvolved in the I/O data path.

Our contributions are:

\begin{enumerate}[leftmargin=*,topsep=2pt,itemsep=1pt]

\item \textbf{First demonstration of GPU-initiated NVMe I/O on consumer
  hardware.} We show that a GeForce RTX 3090 on an AMD B450 platform can
  autonomously issue NVMe commands, achieving up to 3{,}350\MBs{} (78--99\% of
  PCIe~3.0 link bandwidth) without any CPU involvement in the data path.

\item \textbf{Tiered architecture for heterogeneous P2P support.} We propose a
  three-tier approach that degrades gracefully from full P2P (enterprise GPUs)
  to write-only P2P (consumer AMD platforms), requiring no custom kernel
  modules for Tier~1 operation.

\item \textbf{Characterization of PCIe P2P behavior on AMD consumer
  platforms.} We document the asymmetric P2P behavior of AMD B450/Vermeer:
  posted writes succeed through the data fabric while non-posted reads are
  silently dropped, and repeated failed reads cause NVMe link failure requiring
  power cycle recovery.

\item \textbf{Comprehensive evaluation against four baselines.} We compare
  \sysname{} against CPU \texttt{memcpy}, CPU pinned transfer, and NVIDIA
  cuFile (GDS) across block sizes from 4\,KB to 512\,KB and queue depths from
  1 to 32, measuring throughput, latency, IOPS, and CPU utilization.

\item \textbf{Practical LLM inference pipeline.} We demonstrate \sysname{}'s
  applicability to streaming Llama~70B inference on a \$2{,}000 consumer
  system, eliminating the CPU bottleneck in the weight-loading pipeline.

\end{enumerate}

% ============================================================
% 2. BACKGROUND AND MOTIVATION
% ============================================================
\section{Background and Motivation}
\label{sec:background}

\subsection{NVMe Protocol Essentials}

NVMe (Non-Volatile Memory Express)~\cite{nvme-spec} is a host controller
interface for PCIe-attached solid-state storage. The protocol is based on
\emph{submission queues} (SQs) and \emph{completion queues} (CQs) that reside
in host memory:

\begin{itemize}[leftmargin=*,topsep=2pt,itemsep=1pt]
\item The \emph{host} writes 64-byte command entries into the SQ.
\item The host writes the SQ \emph{tail doorbell} register on the controller's
  BAR0 MMIO region to notify the controller of new commands.
\item The \emph{controller} DMAs the command from the SQ, executes it, and
  DMAs a 16-byte completion entry into the CQ.
\item The host detects completions by polling a \emph{phase bit} in the CQ
  entry that toggles on each queue wraparound.
\item The host writes the CQ \emph{head doorbell} to signal consumed
  completions.
\end{itemize}

For data transfers larger than a single page (4\,KB), NVMe uses
\emph{Physical Region Page} (PRP) lists---arrays of physical page addresses
that describe the scatter-gather list for the DMA transfer. For two-page
transfers, PRP2 contains the second page address directly; for larger
transfers, PRP2 points to a page-aligned list of physical addresses.

Critically, NVMe makes no assumption about \emph{who} writes the SQ entries
and doorbell registers. As long as the command format is correct and the
doorbell value is valid, the controller will process the command regardless of
whether a CPU or GPU (or any other bus master) wrote it.

\subsection{PCIe Peer-to-Peer on Consumer Platforms}

PCIe peer-to-peer (P2P) communication allows devices to communicate directly
without routing transactions through system memory. P2P support varies
significantly across platforms:

\textbf{Intel Xeon / AMD EPYC:} Enterprise platforms generally support both
posted (Memory Write) and non-posted (Memory Read) P2P transactions between
PCIe devices.

\textbf{AMD Consumer (B450/B550/X570, Matisse/Vermeer):} The AMD data fabric
forwards \emph{posted} Memory Write TLPs between root ports but \emph{does
not} generate completions for non-posted Memory Read TLPs directed at peer
devices. Read TLPs are silently dropped, causing Completion Timeout (CmpltTO)
errors on the requesting device.

\textbf{NVIDIA GeForce:} NVIDIA's proprietary driver disables P2P on consumer
GPUs via capability checks in \texttt{libcuda.so}. However, the underlying
hardware (GA102/RTX 3090) is physically capable of issuing PCIe posted writes
through its GPU-visible mappings.

This asymmetry---writes work, reads fail---is the foundation of our Tier~1
approach.

\subsection{GPU MMIO and PTX Instructions}

Standard CUDA \texttt{volatile} pointer dereferences do not generate the
correct PCIe transactions for MMIO access. The CUDA compiler may coalesce
or reorder accesses through the GPU's L2 cache hierarchy. To guarantee
correct PCIe BAR register access, we use PTX assembly instructions:

\begin{lstlisting}[language={[x86masm]Assembler},caption={PTX MMIO instructions for PCIe BAR access},label={lst:ptx-mmio}]
// Write: generates PCIe posted MemWr TLP
st.relaxed.mmio.sys.u32 [addr], val;

// Read: generates PCIe non-posted MemRd TLP
ld.relaxed.mmio.sys.u32 val, [addr];
\end{lstlisting}

The \texttt{.mmio} qualifier ensures the access bypasses GPU caches and reaches
the PCIe bus. The \texttt{.sys} qualifier provides system-scope visibility,
ensuring the transaction is visible to all agents in the system (CPU, other
GPUs, PCIe devices).

To register NVMe BAR0 with the GPU's address space, we use
\texttt{cudaHostRegisterIoMemory}, which on our kernel (6.17) required patching
the NVIDIA open-source kernel module to fix a PFN resolution path broken by
the removal of \texttt{follow\_pfn()} in Linux~6.12.

\subsection{Motivation: LLM Inference from Storage}

Consider inference with Llama~2 70B in Q6\_K quantization ($\sim$42\,GB,
669\,MB per layer). With 24\,GB of VRAM, the entire model cannot reside in GPU
memory. A \emph{Streaming Layer Execution Pipeline} (SLEP) reads one layer at a
time from storage, executes the layer's attention and feed-forward computations
on the GPU, then overwrites the layer buffer with the next layer.

The current CPU-mediated pipeline achieves only $\sim$1.5--2\GBs{} effective
throughput due to the double copy (mmap $\rightarrow$ CPU memcpy $\rightarrow$
pinned staging $\rightarrow$ H2D DMA), yielding 0.02 tokens/s. By eliminating
the CPU from the data path, \sysname{} achieves 2.1\GBs{} on a WD SN530
(PCIe~3.0) and 3.35\GBs{} on a WD SN740, directly improving inference
throughput.

% ============================================================
% 3. SYSTEM DESIGN
% ============================================================
\section{System Design}
\label{sec:design}

\subsection{Overview}

\sysname{} enables a CUDA kernel running on a consumer GPU to autonomously
manage NVMe I/O queues, submit read/write commands, and poll completions.
Figure~\ref{fig:architecture} shows the system architecture.

\begin{figure}[t]
\centering
\small
\begin{verbatim}
        GPU (RTX 3090)             NVMe (SN530)
       +-------------+           +-----------+
       | CUDA kernel | ==MMIO==> | BAR0      |
       | (1 thread)  | doorbell  | registers |
       |  sq_submit  | writes    |           |
       |  cq_poll    |           +-----------+
       +------+------+                |
              |                       | DMA
              v                       v
       +------+------+           +-----------+
       | Host pinned |  <==DMA== | NVMe DMA  |
       | memory      |   data    | engine    |
       | SQ/CQ/Data  |           |           |
       +-------------+           +-----------+
\end{verbatim}
\caption{Tier~1 architecture. The GPU writes doorbell registers via MMIO
  posted writes. All queues and data buffers reside in host pinned memory,
  accessible to both the GPU (via PCIe) and the NVMe controller (via DMA).
  The CPU is not involved in the I/O data path.}
\label{fig:architecture}
\end{figure}

The system consists of two phases: a \emph{setup phase} (CPU-driven) and a
\emph{data phase} (GPU-driven).

\textbf{Setup phase (CPU):}
\begin{enumerate}[leftmargin=*,topsep=2pt,itemsep=1pt]
\item Map NVMe BAR0 via \texttt{mmap} of the sysfs resource file.
\item Register BAR0 with the GPU via \texttt{cudaHostRegisterIoMemory}.
\item Initialize the NVMe controller (disable, configure CC, enable, wait
  for CSTS.RDY).
\item Issue admin commands (Identify Controller, Identify Namespace).
\item Create I/O submission and completion queues in host pinned memory.
\item Allocate data buffers and build PRP lists using
  \texttt{/proc/self/pagemap}.
\item Populate a \texttt{gpu\_nvme\_queue} struct with all pointers and pass
  it to the GPU kernel.
\end{enumerate}

\textbf{Data phase (GPU):}
\begin{enumerate}[leftmargin=*,topsep=2pt,itemsep=1pt]
\item Write a 64-byte SQ entry (NVMe Read command) to host pinned memory.
\item Issue \texttt{\_\_threadfence\_system()} to ensure global visibility.
\item Write the SQ tail doorbell to NVMe BAR0 via \texttt{st.relaxed.mmio.sys}.
\item Issue \texttt{\_\_threadfence\_system()} to ensure the doorbell reaches PCIe.
\item The NVMe controller asynchronously DMAs the SQ entry, executes the read,
  DMAs data to the buffer, and writes a CQ entry.
\item Poll the CQ entry's phase bit from host pinned memory.
\item Write the CQ head doorbell to NVMe BAR0.
\item Repeat from step~1 for the next command.
\end{enumerate}

\subsection{Tiered P2P Approach}

We propose three tiers of operation that degrade gracefully based on the
platform's PCIe P2P capabilities:

\begin{description}[leftmargin=*,topsep=4pt,itemsep=2pt]
\item[Tier 1 (Write-Only P2P):] SQ, CQ, and data buffers reside in host
  pinned memory. The GPU writes doorbells to NVMe BAR0 (posted writes) and
  reads SQ/CQ/data from host memory (no P2P reads needed). This tier works on
  AMD consumer platforms where P2P reads fail. \emph{No custom kernel module
  required.}

\item[Tier 2 (Patched P2P):] SQ and CQ remain in host memory; data buffers
  reside in GPU VRAM. The NVMe controller DMAs directly to GPU memory. This
  requires patching the NVIDIA open-source kernel modules to enable P2P on
  GeForce GPUs (similar to the tinygrad P2P patches).

\item[Tier 3 (Native P2P):] All structures (SQ, CQ, data) reside in GPU VRAM.
  This is the BaM~\cite{qureshi2023bam} approach, requiring enterprise GPUs
  (Tesla/A-series) with native P2P DMA support.
\end{description}

\sysname{} currently implements Tier~1, which is sufficient to demonstrate
GPU-autonomous I/O and achieve near-line-rate throughput.

\subsection{Memory Ordering}

Correct operation requires careful memory ordering between the GPU's write
buffer, host DRAM, and the NVMe controller:

\begin{lstlisting}[caption={Critical memory ordering sequence},label={lst:ordering}]
// 1. Write SQ entry to host pinned memory
sqe->opcode = 0x02;  // READ
sqe->prp1 = data_phys;
sqe->cdw10 = lba_low;
// ... (all 16 dwords)

// 2. Fence: ensure SQ writes reach DRAM
__threadfence_system();

// 3. Doorbell MMIO write to NVMe BAR0
mmio_write32(doorbell, sq_tail);

// 4. Fence: ensure doorbell reaches PCIe
__threadfence_system();
\end{lstlisting}

Without the fence between steps 2 and 3, the NVMe controller may process the
doorbell before the SQ entry is visible in host DRAM, causing it to read stale
command data. The \texttt{\_\_threadfence\_system()} instruction provides
system-scope ordering, ensuring all prior writes (to host pinned memory) are
globally visible before the subsequent MMIO write (to NVMe BAR0) is issued.

\subsection{Pipelined I/O}

A single GPU thread can maintain multiple NVMe commands in flight
simultaneously by submitting new commands before polling completions for
previous ones. This is critical for saturating the NVMe device's bandwidth,
as a single outstanding command at 4\,KB cannot keep the SSD busy during its
internal parallelism (flash die interleaving, controller queuing).

The benchmark kernel uses a \emph{pipeline fill} strategy:

\begin{lstlisting}[caption={Pipelined I/O kernel (simplified)},label={lst:pipeline}]
while (completed < num_ops) {
  // Fill pipeline
  while (submitted < num_ops &&
         (submitted - completed) < pipe_depth) {
    sq_submit_read(q, lba[submitted],
                   nlb, prp1[slot], prp2[slot]);
    submitted++;
  }
  // Poll for next completion (any CID)
  cqr = cq_poll_completion(q, timeout);
  if (cqr.success) completed++;
}
\end{lstlisting}

At pipeline depth 32, the GPU thread maintains 32 outstanding NVMe commands.
Each command's SQ entry occupies 64 bytes, and the total SQ footprint is only
2\,KB---negligible relative to the data being transferred.

\subsection{Queue Depth Analysis via Little's Law}
\label{sec:pipeline-analysis}

The effectiveness of our single-thread pipelined design can be understood through
Little's Law: to saturate a device with bandwidth $B$ and per-command service
time $t_s$, the required queue depth is $N = B \cdot t_s / s$, where $s$ is the
block size. For the SN530 at 256\,KB blocks with $B = 3{,}400\MBs$ and
$t_s \approx 75\us$ (measured single-command latency), the minimum queue depth is
$N = 3{,}400 \times 75 / (256 \times 1{,}000) \approx 1.0$. At 4\,KB blocks,
$N = 3{,}400 \times 10.6 / (4 \times 1{,}000) \approx 9.0$.

This analysis explains three observed phenomena. First, \sysname{} at QD=32
provides ample headroom for both block sizes, explaining why we achieve 78\% of
link bandwidth. Second, CPU methods show no queue depth scaling
(Table~\ref{tab:scaling}) because the kernel I/O stack serializes requests---the
effective QD at the device is always $\sim$1 regardless of the application-level
queue depth. Third, a \emph{single} GPU thread suffices because the pipeline
management overhead (SQ entry construction + doorbell write + CQ poll) takes
$\sim$2\us{} per command, far below the NVMe service time, leaving the thread
idle in the CQ polling loop most of the time. Additional threads would provide
negligible benefit unless the NVMe device's internal parallelism required higher
queue depth (e.g., Gen5 devices with sub-microsecond latency).

\subsection{PRP List Construction}

For transfers larger than 4\,KB (one page), the NVMe controller needs
physical addresses for each page of the destination buffer. We pre-build PRP
(Physical Region Page) lists on the CPU during the setup phase:

\begin{enumerate}[leftmargin=*,topsep=2pt,itemsep=1pt]
\item Allocate a contiguous virtual buffer via \texttt{posix\_memalign} +
  \texttt{mlock}.
\item Register with CUDA via \texttt{cudaHostRegister} for GPU accessibility.
\item Open \texttt{/proc/self/pagemap} and translate each virtual page to its
  physical frame number (PFN).
\item Build per-slot PRP lists: PRP1 = first page physical address; for
  2-page transfers, PRP2 = second page address; for $>$2 pages, PRP2 = physical
  address of a page-aligned list of physical addresses.
\item Pass PRP1/PRP2 arrays to the GPU kernel as parameters.
\end{enumerate}

This design amortizes the costly pagemap translation over the entire benchmark
run, keeping the GPU data-path overhead minimal.

% ============================================================
% 4. IMPLEMENTATION
% ============================================================
\section{Implementation}
\label{sec:impl}

\sysname{} is implemented in $\sim$4{,}500 lines of C and CUDA code (excluding
benchmarks and tests), organized into GPU-side device code and CPU-side host
code.

\subsection{GPU-Side Components}

\textbf{mmio\_ops.cuh}: Provides \texttt{mmio\_write32},
\texttt{mmio\_read32}, and their 64-bit variants using PTX inline assembly.
Under simulation (\texttt{GPUNVME\_USE\_SIM}), falls back to volatile pointer
dereferences.

\textbf{sq\_submit.cuh}: Constructs 64-byte NVMe Read/Write command entries
directly in the submission queue. Each field is written as an aligned 32-bit
or 64-bit word to avoid misaligned access penalties. CDW0 (containing the
opcode and command ID) is written as a single atomic 32-bit store.

\textbf{cq\_poll.cuh}: Implements phase-bit-based completion detection. Reads
CQ DW3 (offset~12, naturally aligned) as a single 32-bit word, extracts the
phase bit and command ID via bit shifts. Advances the CQ head pointer with
wraparound and phase flip.

\textbf{queue\_state.cuh}: Defines the \texttt{gpu\_nvme\_queue} struct that
encapsulates all state needed by a GPU kernel to manage an NVMe I/O queue pair
independently: SQ/CQ pointers, doorbell register pointers (BAR0 MMIO),
queue dimensions, current tail/head positions, phase bit, and namespace
information.

\subsection{Host-Side Components}

\textbf{controller.c}: Full NVMe controller initialization per the NVMe
specification: disable controller (CC.EN=0), wait CSTS.RDY=0, configure admin
queues (AQA, ASQ, ACQ), set CC (command set, page size, queue entry sizes),
enable (CC.EN=1), wait CSTS.RDY=1, issue Identify Controller and Identify
Namespace admin commands.

\textbf{io\_queue.c}: Creates I/O queue pairs via admin commands (Create I/O CQ,
Create I/O SQ). Allocates queue memory using \texttt{posix\_memalign} with
\texttt{mlock} to guarantee page alignment (discovered that
\texttt{cudaMallocHost}'s suballocator may return sub-page offsets, violating
NVMe's page-alignment requirement for queue base addresses).

\textbf{dma\_alloc.c}: Resolves virtual-to-physical address translation via
\texttt{/proc/self/pagemap}. Provides PRP list allocation and construction
for multi-page transfers.

\textbf{bar\_map.c}: Maps NVMe BAR0 via \texttt{mmap} of the sysfs
\texttt{resource0} file with \texttt{O\_SYNC} to ensure uncacheable access.
Registers the mapping with the GPU via \texttt{cudaHostRegisterIoMemory |
cudaHostRegisterMapped}.

\subsection{NVIDIA Driver Patch}

On Linux~6.17, \texttt{cudaHostRegisterIoMemory} fails because the NVIDIA
open-source kernel module's PFN resolution path calls
\texttt{nv\_follow\_pfn()}, which delegates to \texttt{follow\_pfn()}---a
function removed in Linux~6.12 (commit \texttt{233eb0bf3b94}). The fallback
\texttt{nv\_follow\_flavors()} unconditionally returns $-1$.

We patched \texttt{os-mlock.c} to compute the PFN from the VMA's
\texttt{vm\_pgoff} field for \texttt{VM\_PFNMAP} mappings:

\begin{lstlisting}[caption={NVIDIA driver PFN resolution patch},label={lst:patch}]
if (vma && (vma->vm_flags & VM_PFNMAP)) {
    *pfn = vma->vm_pgoff +
      ((address - vma->vm_start) >> PAGE_SHIFT);
    return 0;
}
\end{lstlisting}

This works because \texttt{io\_remap\_pfn\_range()} (used by PCI BAR mmap)
stores the starting PFN in \texttt{vm\_pgoff} and creates contiguous mappings.
The GPL-only \texttt{follow\_pfnmap\_start()} API cannot be used by the MIT-licensed
NVIDIA module.

\subsection{VFIO Device Setup}

The test NVMe (WD SN530) is passed through to userspace via VFIO in
no-IOMMU mode. After unbinding from the kernel NVMe driver and binding to
\texttt{vfio-pci}, the device falls to D3 power state. We restore it to D0
and re-enable Memory Space and Bus Master via \texttt{setpci}:

\begin{lstlisting}[language=bash,caption={VFIO setup and power state recovery}]
echo on > /sys/bus/pci/devices/$BDF/power/control
setpci -s $BDF 0x84.W=0x0008  # Force D0
setpci -s $BDF COMMAND=0x0006 # Mem+BusMaster
\end{lstlisting}

% ============================================================
% 5. EVALUATION
% ============================================================
\section{Evaluation}
\label{sec:eval}

\subsection{Experimental Setup}

Table~\ref{tab:hardware} summarizes our hardware configuration.

\begin{table}[t]
\centering
\caption{Hardware configuration}
\label{tab:hardware}
\small
\begin{tabular}{ll}
\toprule
\textbf{Component} & \textbf{Detail} \\
\midrule
GPU & NVIDIA RTX 3090 (GA102, 24\,GB, PCIe 3.0 $\times$8\textsuperscript{$\dagger$}) \\
CPU & AMD Ryzen 7 5800X (8C/16T, 3.8--4.7\,GHz) \\
Platform & AMD B450, DDR4-3200 (48\,GB) \\
NVMe (test 1) & WD SN530 1\,TB, PCIe 3.0 $\times$4 (VFIO) \\
NVMe (test 2) & WD SN740 512\,GB, PCIe 3.0 $\times$4 (VFIO) \\
NVMe (system) & Samsung 980 PRO 500\,GB, PCIe 4.0 $\times$4 \\
OS & Ubuntu 25.10 (kernel 6.17.0) \\
CUDA & 13.1, Driver 590.48.01 (open, patched) \\
\bottomrule
\end{tabular}

\smallskip
\noindent\textsuperscript{$\dagger$}B450 shares PCIe lanes between the primary
$\times$16 slot and M.2\_1: when an M.2 NVMe is populated, the GPU
negotiates PCIe~3.0~$\times$8 ($\sim$6.5\GBs{} host-to-device bandwidth).
The SN740 is a PCIe~4.0 device but operates at Gen3 speeds on this platform.
\end{table}

\textbf{Note on device asymmetry.} The detailed per-block-size benchmarks
(Tables~\ref{tab:throughput}--\ref{tab:scaling}) use the WD SN530
(PCIe~3.0~$\times$4, $\sim$3{,}400\MBs{} theoretical max) as the VFIO
device. The CPU baselines (memcpy, pinned, cuFile) use the Samsung 980 PRO
(PCIe~4.0~$\times$4, $\sim$7{,}000\MBs{} theoretical max) because the SN530 has no
kernel device path. Despite this 2$\times$ asymmetry favoring the CPU baselines,
\sysname{} outperforms them due to its ability to pipeline I/O operations.
We additionally validate on a WD SN740 (PCIe~4.0 device, Gen3 link on B450),
which achieves 3{,}350\MBs{} sustained---near the Gen3~$\times$4 theoretical
maximum of $\sim$3{,}400\MBs{}.

To enable fair comparison, we normalize throughput as a fraction of each device's
theoretical PCIe link bandwidth. At 256\,KB/QD=32, \sysname{} achieves
2{,}666/3{,}400 = \textbf{78.4\%} of its Gen3 link on the SN530 and
3{,}350/3{,}400 = \textbf{98.5\%} on the SN740, while \texttt{cpu\_memcpy}
achieves 1{,}731/7{,}000 = \textbf{24.7\%}, \texttt{cpu\_pinned} reaches
1{,}409/7{,}000 = 20.1\%, and cuFile 1{,}396/7{,}000 = 19.9\%. The GPU-initiated
path thus extracts 3--4$\times$ more of the available link bandwidth than
CPU-mediated methods. This gap is explained by Little's Law
(\S\ref{sec:pipeline-analysis}): saturating NVMe bandwidth requires
$\text{QD} = \text{BW} \times \text{latency}$, and only \sysname{} maintains
sufficient queue depth from the device's perspective.

\textbf{Methodology:} Each configuration (block size $\times$ queue depth) is
run 3 times. We measure per-completion latency using the GPU clock
(\texttt{clock64()}), wall-clock time, and CPU utilization via
\texttt{/proc/stat}. All tests use sequential access patterns. Block sizes
range from 4\,KB to 512\,KB; queue depths from 1 to 32.

\subsection{Throughput}

Figure~\ref{fig:throughput-qd} shows throughput scaling with queue depth for
each method.

\begin{table}[t]
\centering
\caption{Throughput (MB/s) at QD=32 by block size. \sysname{} uses the
  slower SN530; CPU baselines use the faster 980~PRO. Despite this, \sysname{}
  achieves the highest throughput at all block sizes.}
\label{tab:throughput}
\small
\begin{tabular}{lrrrr}
\toprule
\textbf{Block Size} & \textbf{gpu\_direct} & \textbf{cpu\_memcpy} & \textbf{cpu\_pinned} & \textbf{cuFile} \\
\midrule
4\,KB   & \textbf{336}   & 180  & 153  & 151 \\
16\,KB  & \textbf{1{,}393} & 626  & 516  & 510 \\
64\,KB  & \textbf{2{,}111} & 1{,}373 & 1{,}152 & 1{,}144 \\
256\,KB & \textbf{2{,}666} & 1{,}731 & 1{,}409 & 1{,}396 \\
512\,KB & \textbf{2{,}634} & 2{,}008 & 1{,}700 & 1{,}696 \\
\bottomrule
\end{tabular}
\end{table}

\textbf{Key finding 1: Queue depth scaling.} \sysname{} shows dramatic
throughput improvement with queue depth: 2.0$\times$ at 4\,KB (173 to
336\MBs{}) and 2.3$\times$ at 16\,KB (615 to 1{,}393\MBs{}). In contrast,
CPU-mediated methods show \emph{no meaningful queue depth scaling}---their
throughput is flat across QD=1 to QD=32 because the CPU serializes I/O
operations through \texttt{memcpy}.

\textbf{Key finding 2: Link bandwidth utilization.} At 256\,KB/QD=32, \sysname{}
achieves 2{,}666\MBs{} on the SN530 (78\% of Gen3~$\times$4 link) and
3{,}350\MBs{} on the SN740 (99\% of Gen3~$\times$4 link). The normalized
comparison above shows that CPU methods utilize only 20--25\% of their faster
Gen4 link, confirming that the GPU-initiated path is fundamentally more
efficient at extracting device bandwidth.

\textbf{Key finding 3: Advantage despite slower hardware.} Even though the
CPU baselines use a 2$\times$ faster SSD, \sysname{} achieves 1.3--2.2$\times$
higher throughput across all block sizes at QD=32. Per the Little's Law analysis
(\S\ref{sec:pipeline-analysis}), this is because the CPU-mediated path
effectively operates at QD=1 from the device's perspective, regardless of the
application-level queue depth.

\subsection{Queue Depth Scaling Factor}

Table~\ref{tab:scaling} quantifies the scaling advantage. The GPU's ability to
maintain 32 in-flight commands from a single thread is the fundamental
differentiator.

\begin{table}[t]
\centering
\caption{Throughput scaling factor (QD=32 / QD=1)}
\label{tab:scaling}
\small
\begin{tabular}{lrrrr}
\toprule
\textbf{Block Size} & \textbf{gpu\_direct} & \textbf{cpu\_memcpy} & \textbf{cpu\_pinned} & \textbf{cuFile} \\
\midrule
4\,KB   & 1.95$\times$ & 1.08$\times$ & 0.96$\times$ & 1.01$\times$ \\
16\,KB  & 2.27$\times$ & 1.00$\times$ & 0.98$\times$ & 1.00$\times$ \\
64\,KB  & 1.41$\times$ & 1.10$\times$ & 0.98$\times$ & 1.00$\times$ \\
256\,KB & 1.65$\times$ & 1.00$\times$ & 0.99$\times$ & 1.01$\times$ \\
512\,KB & 1.98$\times$ & 1.02$\times$ & 1.00$\times$ & 0.99$\times$ \\
\bottomrule
\end{tabular}
\end{table}

\subsection{Latency}

Table~\ref{tab:latency} shows median per-operation latency. At QD=1, each
measurement represents the full NVMe round-trip; at QD=32, it measures the
inter-completion gap (time between successive completions arriving at the CQ).

\begin{table}[t]
\centering
\caption{Median latency (\us{}) for \sysname{}}
\label{tab:latency}
\small
\begin{tabular}{lrrrr}
\toprule
\textbf{Block Size} & \textbf{QD=1} & \textbf{QD=4} & \textbf{QD=16} & \textbf{QD=32} \\
\midrule
4\,KB   & 10.6 & 1.70 & 1.49 & 1.50 \\
16\,KB  & 14.2 & 15.5 & 1.48 & 1.47 \\
64\,KB  & 29.0 & 33.4 & 13.0 & 14.0 \\
256\,KB & 126  & 69   & 63   & 63 \\
512\,KB & 412  & 154  & 141  & 141 \\
\bottomrule
\end{tabular}
\end{table}

At 4\,KB/QD=32, the inter-completion latency is only 1.5\us{}, demonstrating
that the GPU's CQ polling loop is extremely efficient. The latency standard
deviation at this configuration is 0.23\us{}, indicating highly consistent
completion timing.

At QD=1, \sysname{} achieves lower latency than CPU methods for small blocks
(10.6\us{} vs. 19.6\us{} for \texttt{cpu\_memcpy} at 4\,KB), a 46\% reduction.
This is because the GPU writes the doorbell directly via MMIO, avoiding the
CPU's system call, I/O scheduler, and interrupt handling overhead.

\subsection{CPU Utilization}

\begin{table}[t]
\centering
\caption{CPU utilization (\%) across all configurations}
\label{tab:cpu}
\small
\begin{tabular}{lrrr}
\toprule
\textbf{Method} & \textbf{Mean} & \textbf{Min} & \textbf{Max} \\
\midrule
gpu\_direct  & 8.5  & 5.6  & 15.0 \\
cpu\_memcpy  & 4.3  & 0.0  & 12.2 \\
cpu\_pinned  & 4.7  & 0.0  & 12.5 \\
cuFile       & 5.1  & 2.1  & 10.4 \\
\bottomrule
\end{tabular}
\end{table}

The CPU utilization for \sysname{} (8.5\% average) reflects CUDA runtime
overhead (kernel launch, synchronization), \emph{not} I/O data path work. The
CPU is entirely free to perform other tasks (inference computation,
preprocessing) during GPU-driven I/O. In contrast, the CPU methods block a
thread on I/O completion, making that thread unavailable for computation
despite showing lower aggregate CPU utilization across all 16 logical cores.

\subsection{cuFile on Consumer Hardware}

A notable finding is that NVIDIA cuFile (GPUDirect Storage API) performs
identically to \texttt{cpu\_pinned} on our GeForce RTX 3090 hardware. Across
all 60 configurations, the two methods differ by less than 3\% on average.
This confirms that GDS falls back to a CPU-mediated path on consumer GPUs,
providing no hardware acceleration. \sysname{} fills this gap by providing
genuine GPU-initiated I/O on hardware that GDS does not support.

\subsection{Large Sequential Reads}

For the LLM inference use case, we measured throughput for large sequential
reads representative of loading model layers:

\begin{table}[t]
\centering
\caption{Large sequential read throughput (QD=32, 512\,KB blocks)}
\label{tab:large-read}
\small
\begin{tabular}{lrr}
\toprule
\textbf{Transfer Size} & \textbf{SN530 (MB/s)} & \textbf{SN740 (MB/s)} \\
\midrule
4\,MB    & 2{,}122 & --- \\
128\,MB  & 2{,}734 (peak) & --- \\
669\,MB (1 layer) & 2{,}127 & 3{,}350 \\
8.6\,GB sustained & --- & 3{,}350 \\
\bottomrule
\end{tabular}
\end{table}

The SN530 peaks at 2{,}734\MBs{} at 128\,MB, dropping to 2{,}127\MBs{} at
669\,MB due to thermal throttling during sustained reads. The SN740 sustains
3{,}350\MBs{} even at 8.6\,GB---near the Gen3~$\times$4 theoretical
maximum---demonstrating that \sysname{} can fully saturate the PCIe link for
layer-sized transfers.

\subsection{LLM Inference Performance}

To contextualize the I/O throughput results, we present both \emph{projected}
token rates from raw I/O bandwidth and \emph{measured} results from
integrating \sysname{} into a streaming inference engine (ntransformer) for
Llama~3.1-70B in Q6\_K quantization ($\sim$42\,GB, 80~layers,
$\sim$669\,MB/layer).

Table~\ref{tab:llm-projection} shows projected rates for the pure I/O-bound
regime where each token requires streaming the entire model from NVMe
(no caching).

\begin{table}[t]
\centering
\caption{Llama-70B inference: projected (I/O-bound, full model per token)
  and measured (integrated with tiered caching)}
\label{tab:llm-projection}
\small
\begin{tabular}{lrr}
\toprule
\textbf{Data Path} & \textbf{BW} & \textbf{tok/s} \\
\midrule
\multicolumn{3}{l}{\emph{Projected (80 layers from NVMe per token)}} \\
mmap+memcpy                   & 1.5\GBs{} & 0.028 \\
\sysname{} (SN530)            & 2.1\GBs{} & 0.039 \\
\sysname{} (SN740)            & 3.35\GBs{} & 0.063 \\
\midrule
\multicolumn{3}{l}{\emph{Measured (ntransformer integration)}} \\
\sysname{} NVMe-only (SN740)  & 3.35\GBs{} & 0.06 \\
\sysname{} NVMe-only (SN530)  & 2.1\GBs{} & 0.04 \\
Tiered VRAM/RAM (no NVMe)     & 6.5\GBs{} & 0.20 \\
Tiered + layer skip (0.98)    & 6.5\GBs{} & 0.27 \\
Tiered + Q4\_K\_M + skip      & 6.5\GBs{} & \textbf{0.50} \\
\bottomrule
\end{tabular}
\end{table}

\textbf{Measured integration results.} We integrated \sysname{} into
ntransformer, a streaming LLM inference engine with tiered caching
(VRAM/RAM/NVMe). On the SN740, \sysname{} achieves 0.06~tok/s for
70B Q6\_K with all 80 layers streamed from NVMe per token---a 2$\times$
improvement over mmap+memcpy (0.03~tok/s) and consistent with the
projected I/O bandwidth. The CPU is entirely free during NVMe transfers,
enabling concurrent compute overlap.

When combined with tiered caching (29 layers in VRAM, 51 in RAM, 0 on
NVMe), the system achieves 0.20~tok/s. The bottleneck shifts from NVMe
bandwidth to PCIe H2D transfers at Gen3~$\times$8 ($\sim$6.5\GBs{},
103\,ms per 669\,MB layer). Adding adaptive layer skipping
(threshold~0.98, skipping 16--20 middle layers) yields 0.27~tok/s, and
combining with Q4\_K\_M quantization reaches 0.50~tok/s---a
17$\times$ improvement over the CPU-mediated NVMe baseline.

\textbf{Important context.} The absolute token rates for NVMe-only streaming
are inherently low because the entire model must transit from storage per
token. For reference, llama.cpp achieves $\sim$1~tok/s for 70B models when
weights reside in GPU VRAM across two GPUs (tensor parallelism), and
FlexGen~\cite{sheng2023flexgen} achieves 1~tok/s for OPT-175B with
aggressive CPU/disk offloading. The contribution of \sysname{} is not the
absolute token rate but the \textbf{removal of the CPU from the storage data
path}: measured CPU utilization during NVMe-direct transfers is $\sim$0\%
versus 100\% (one core saturated) for mmap+memcpy.

For Q8\_0 quantization ($\sim$70\,GB) that \emph{exceeds} our 48\,GB RAM,
NVMe streaming via \sysname{} is the only viable path without adding
system memory. Combined with NanoQuant~\cite{chong2026nanoquant} (sub-1-bit
quantization compressing 70B to $\sim$8\,GB), the data volume per token
drops from 42\,GB to $\sim$8\,GB, projecting $\sim$0.5~tok/s from NVMe
alone---matching the tiered+optimized result but without requiring RAM
caching.

% ============================================================
% 6. BUGS AND ENGINEERING CHALLENGES
% ============================================================
\section{Engineering Challenges}
\label{sec:bugs}

Developing \sysname{} required solving several non-trivial challenges at the
intersection of GPU programming, NVMe protocol, PCIe topology, and kernel
driver internals. We document these for reproducibility.

\begin{table}[t]
\centering
\caption{Critical bugs discovered and resolved}
\label{tab:bugs}
\small
\begin{tabular}{clp{0.55\textwidth}}
\toprule
\textbf{\#} & \textbf{Bug} & \textbf{Root Cause \& Fix} \\
\midrule
1 & \texttt{cudaHostRegisterIoMemory} fails &
    \texttt{follow\_pfn()} removed in kernel~6.12. Patched
    \texttt{os-mlock.c}: compute PFN from \texttt{vm\_pgoff}. \\
2 & GPU reads return \texttt{0xFFFFFFFF} &
    AMD root complex drops non-posted P2P TLPs. Use Tier~1 (writes only). \\
3 & Admin commands hang &
    Admin CQ not page-aligned (\texttt{cudaMallocHost} suballocator).
    Allocate $\geq$4096 bytes. \\
4 & Pipeline depth $\geq$4 timeouts &
    NVMe completions arrive out-of-order; poll discarded non-matching CQEs.
    Use \texttt{cq\_poll\_completion} (accept any CID). \\
5 & Post-warmup timeouts &
    Resetting queue state while NVMe controller retains
    internal pointers. Leave queue state rolling. \\
6 & GPU reads crash NVMe link &
    Accumulated CmpltTO errors from failed P2P reads. Avoid GPU reads;
    power cycle to recover. \\
7 & 64-bit MMIO write corruption &
    Non-atomic 64-bit write to BAR0. Split into two 32-bit writes
    (low word first per NVMe spec). \\
\bottomrule
\end{tabular}
\end{table}

Bug~\#4 deserves special attention. NVMe controllers may reorder completions
within a queue---returning CQE for command $B$ before command $A$ even if $A$
was submitted first. Our initial implementation used
\texttt{cq\_poll\_for\_cid}, which discarded CQEs not matching the expected
command ID. When the pipeline had multiple in-flight commands, this caused
valid completions to be lost, leading to timeouts. The fix was to poll for
\emph{any} completion (\texttt{cq\_poll\_completion}) and track which command
IDs have completed, accepting completions in whatever order the NVMe
controller delivers them.

% ============================================================
% 7. RELATED WORK
% ============================================================
\section{Related Work}
\label{sec:related}

\textbf{GPU-Initiated Storage Access.}
BaM~\cite{qureshi2023bam} (ASPLOS~'23) is the closest prior work, demonstrating
GPU-orchestrated on-demand storage access. BaM provisions NVMe queues directly
in GPU HBM via GPUDirect RDMA and uses GPUDirect Async for doorbell writes,
achieving 45.8M~IOPs across 10~SSDs. GIDS~\cite{park2024gids} (VLDB~'24)
extends BaM to GNN training with a dynamic storage access accumulator and
window buffering, reaching 582$\times$ speedup over DGL on terabyte-scale
graphs. Both systems require enterprise GPUs (A100) with native P2P DMA
support, PCIe expansion chassis, and custom kernel modules. \sysname{} differs
fundamentally: it works on consumer hardware (\$2{,}000 vs.\ \$10{,}000+)
where full P2P is unavailable, using only PCIe posted writes for doorbell
access and host pinned memory for data staging.

GPUfs~\cite{silberstein2014gpufs} pioneered GPU-initiated file I/O but routes
all I/O through a CPU-side daemon. DRAGON~\cite{markthub2018dragon} extends GPU
memory to NVM via page faults, introducing high per-fault latency.
SPIN~\cite{bergman2017spin} integrates P2P DMA into the Linux page cache
but requires PCIe switch topology and CPU-initiated I/O.
Phoenix~\cite{phoenix2025} refactors the GDS I/O stack, mapping GPU memory into
PCIe BAR space for direct NVMe access---but remains limited to enterprise GPUs.

\textbf{GPUDirect Storage.}
NVIDIA's GDS~\cite{gds} enables direct NVMe-to-GPU DMA; as of CUDA~12.8, it
supports P2P DMA via the upstream kernel PCI P2PDMA infrastructure on x86\_64.
However, GDS requires enterprise GPUs, certified storage, and the
\texttt{nvidia-fs} kernel module. Our evaluation confirms that cuFile provides
no benefit on GeForce hardware, falling back to a CPU-mediated path identical
to \texttt{cpu\_pinned}.

\textbf{Userspace NVMe Drivers.}
SPDK~\cite{spdk} provides high-performance CPU-driven userspace NVMe access via
polling. Our work extends the userspace NVMe concept to the GPU: instead of CPU
cores polling SQ/CQ, a CUDA thread does. libnvm~\cite{markussen2021smartio}
enables GPU DMA from NVMe but requires a kernel module and P2P-capable platforms.

\textbf{LLM Inference on Constrained Hardware.}
A rapidly growing body of work addresses LLM inference with limited GPU memory,
all using CPU-mediated storage I/O:
FlexGen~\cite{sheng2023flexgen} (ICML~'23) offloads weights, activations, and
KV cache across GPU/CPU/disk using linear programming, achieving OPT-175B on a
single 16\,GB GPU at 1~tok/s.
PowerInfer~\cite{song2023powerinfer} (SOSP~'24) exploits activation sparsity
(power-law neuron distribution) to preload hot neurons on GPU and compute cold
neurons on CPU, reaching 11.69$\times$ over llama.cpp on RTX~4090.
PIPO~\cite{liu2025pipo} designs a fine-grained offloading pipeline for consumer
devices, increasing GPU utilization from $<$40\% to $>$90\% on RTX~3060.
HeadInfer~\cite{luo2025headinfer} offloads KV cache to CPU RAM at per-attention-head
granularity, enabling 4M-token inference with Llama~3-8B on a single RTX~4090.
ServerlessLLM~\cite{fu2024serverlessllm} (OSDI~'24) exploits the full
near-GPU storage hierarchy for fast checkpoint loading, achieving 12\GBs{}
from RAID0-NVMe with optimized \texttt{O\_DIRECT} and pinned memory pipelining.
InstInfer~\cite{instinfer2024} offloads attention computation to computational
storage devices. DeepSpeed-Inference~\cite{aminabadi2022deepspeed} supports
NVMe offloading via ZeRO-Infinity. LLM in a Flash~\cite{alizadeh2023llmflash}
targets Apple's unified memory. NanoQuant~\cite{chong2026nanoquant} compresses
LLMs to sub-1-bit levels, enabling Llama~2-70B on an 8\,GB consumer GPU---a
technique complementary to \sysname{} that would dramatically reduce the data
volume needing to be streamed from NVMe.

None of these systems achieve \emph{GPU-initiated} storage access; all route
I/O through the CPU. A recent I/O characterization study~\cite{cheops2025io}
of LLM offloading to NVMe confirms that current frameworks do not saturate
SSD bandwidth---model offloading is dominated by 128\,KiB sequential reads
at well below device peak. \sysname{} addresses this gap by removing the CPU
from the storage data path.

\textbf{PCIe P2P and Interconnect Evolution.}
Prior work on PCIe P2P~\cite{markussen2021smartio} has characterized P2P
behavior across Intel and some AMD platforms. The Linux P2PDMA subsystem
provides growing kernel support for device-to-device DMA. Looking forward,
CXL memory pools~\cite{zhong2025cxl} could provide shared memory regions
accessible by both GPU and NVMe, simplifying the data path. However, none
of this prior work has documented the specific asymmetry on AMD B450/Vermeer
(posted writes succeed, non-posted reads cause link failure) that we
characterize and exploit in \sysname{}.

\subsection{Positioning}

Table~\ref{tab:comparison} positions \sysname{} against related systems.
\sysname{} is the only system that combines GPU-initiated I/O with consumer
hardware compatibility.

\begin{table}[t]
\centering
\caption{Feature comparison with related systems}
\label{tab:comparison}
\small
\begin{tabular}{lccccc}
\toprule
\textbf{System} & \rotatebox{60}{\textbf{GPU initiates}} & \rotatebox{60}{\textbf{No CPU data}} & \rotatebox{60}{\textbf{Consumer HW}} & \rotatebox{60}{\textbf{No kmod}} & \rotatebox{60}{\textbf{No ent. GPU}} \\
\midrule
\sysname{}       & \checkmark & \checkmark & \checkmark & \checkmark & \checkmark \\
BaM              & \checkmark & \checkmark & --         & --         & -- \\
GIDS             & \checkmark & \checkmark & --         & --         & -- \\
GDS              & --         & \checkmark & --         & --         & -- \\
SPDK             & --         & \checkmark & \checkmark & --         & N/A \\
ServerlessLLM    & --         & --         & --         & \checkmark & \checkmark \\
PowerInfer       & --         & --         & \checkmark & \checkmark & \checkmark \\
FlexGen          & --         & --         & \checkmark & \checkmark & \checkmark \\
PIPO             & --         & --         & \checkmark & \checkmark & \checkmark \\
\bottomrule
\end{tabular}
\end{table}

% ============================================================
% 8. DISCUSSION AND FUTURE WORK
% ============================================================
\section{Discussion and Future Work}
\label{sec:discussion}

\textbf{Tier~1 bandwidth ceiling.} Tier~1 routes data through host pinned
memory, adding one PCIe hop compared to direct NVMe-to-GPU DMA. On our B450
platform, the GPU's PCIe~3.0~$\times$8 link to host memory provides
$\sim$6.5\GBs{} of bandwidth---roughly 2$\times$ the SN530's 3.4\GBs{}.
The NVMe device is therefore the bottleneck, not the extra hop, and Tier~1
loses no throughput from the indirection. However, the margin is tight: a
Gen4~$\times$4 NVMe at full speed (7\GBs{}) would exceed our GPU link
bandwidth, and even on our platform the SN740 reaches 3.35\GBs{} ($\sim$52\%
of the GPU link). On platforms with full Gen4~$\times$16 (25\GBs{}), this
headroom increases substantially. For our B450 configuration, this motivates
\textbf{Tier~2}: enabling NVMe DMA directly to GPU VRAM via the tinygrad P2P
patches for the NVIDIA open-source kernel module, which would eliminate the
host memory intermediate and become essential when upgrading to faster NVMe
devices or adding a second drive.

\textbf{Why a single GPU thread suffices.} Our Little's Law analysis
(\S\ref{sec:pipeline-analysis}) shows that QD=32 provides ample headroom to
saturate the SN530. The I/O management thread spends $\sim$2\us{} per command
(SQ construction + doorbell write + fence), far below the device service time,
and resides in a tight CQ polling loop otherwise. This consumes exactly one
warp slot, leaving the remaining 81 SMs (on GA102) free for tensor computation.
A multi-warp design managing multiple NVMe queues could provide higher
aggregate throughput for multi-device setups, but for a single NVMe device, one
thread is sufficient. We plan to evaluate multi-queue parallelism with 2--4
NVMe devices to determine scaling behavior.

\textbf{Bandwidth utilization in context.} \sysname{} achieves 78\% of the
PCIe~3.0~$\times$4 link bandwidth on the SN530 and 99\% on the SN740 at
256\,KB/QD=32. For comparison, BaM~\cite{qureshi2023bam} reports 90\% of link
bandwidth on enterprise A100 hardware with GPUDirect RDMA and queues in GPU
HBM. The SN530's 78\% gap is attributable to Tier~1's extra host memory hop
and device-side throttling; the SN740's 99\% demonstrates that \sysname{} can
match Tier~3-level link utilization on consumer hardware when the drive sustains
its throughput. Despite the extra hop, the normalized comparison
(\S\ref{sec:eval}) shows that \sysname{} extracts 3--4$\times$ more of the
available link bandwidth than CPU-mediated methods on the same platform.

\textbf{Synergy with quantization.} The measured 0.50~tok/s result
(Table~\ref{tab:llm-projection}) already demonstrates this synergy:
Q4\_K\_M quantization reduces per-layer size from 669\,MB to $\sim$370\,MB,
and combined with layer skipping, the effective data volume per token drops
substantially. NanoQuant~\cite{chong2026nanoquant} pushes further with
sub-1-bit quantization compressing Llama~2-70B to $\sim$8\,GB total.
Combining \sysname{} NVMe-only streaming with NanoQuant would reduce the
per-token transfer from 42\,GB to $\sim$8\,GB, projecting $\sim$0.5~tok/s
from NVMe alone---matching the tiered+optimized result but without requiring
RAM caching. This synergy is bidirectional: quantization reduces I/O volume
while GPU-initiated I/O reduces transfer latency, and neither technique
substitutes for the other.

\textbf{CXL and future interconnects.} CXL~3.0 memory pools~\cite{zhong2025cxl}
could fundamentally change the P2P landscape by providing cache-coherent
shared memory regions accessible to both GPU and NVMe controllers. In such
architectures, \sysname{}'s Tier~1 host-memory approach would naturally
benefit from the lower-latency CXL path. More importantly, CXL-attached
NVMe devices could enable Tier~3-equivalent semantics on consumer platforms
if the CXL fabric supports GPU-initiated transactions.

\textbf{Safety and reliability.} Operating NVMe at the register level bypasses
all kernel-level safety mechanisms. Bugs in command construction can cause
data corruption or device hangs (Bug~\#6, Table~\ref{tab:bugs}). We mitigate
this through a simulator mode (verifying logic without hardware),
comprehensive struct layout tests, and careful power management. In
production, a watchdog timer in the GPU kernel detects hangs (via
\texttt{clock64()}) and returns error codes. A formal verification of the
command construction path would further strengthen reliability.

\textbf{Portability.} \sysname{}'s Tier~1 approach should work on any platform
where the GPU can issue PCIe posted writes to a peer device's BAR. This
includes AMD EPYC (full P2P), Intel Xeon (full P2P), and AMD consumer (posted
writes only). Platforms with active PCIe switches (e.g., PLX/Broadcom) should
also support the necessary posted write routing. We plan to evaluate on Intel
12th/13th-gen consumer platforms, where full P2P may enable Tier~2/3 without
enterprise GPUs.

\textbf{End-to-end integration.} We integrated \sysname{} into ntransformer,
a streaming inference engine with tiered caching (\S\ref{sec:eval}). The
integration confirms that GPU-initiated I/O composes well with a
double-buffering strategy: while the GPU computes attention and feed-forward
on layer $n$, \sysname{} prefetches layer $n+1$ from NVMe. The measured
0.06~tok/s for NVMe-only streaming matches the projected I/O bandwidth,
validating that the primitive performs as expected in an application context.
Further optimization (tiered caching, layer skipping, aggressive quantization)
reaches 0.50~tok/s, demonstrating that \sysname{} serves as a viable building
block for practical single-GPU inference of 70B-class models.

\textbf{Limitations.} (1)~The CPU baselines use a different (faster) NVMe
device than \sysname{}. The bandwidth normalization (\S\ref{sec:eval})
controls for this, and the SN740 results (99\% link utilization) confirm
\sysname{} is not device-specific. (2)~Write support is implemented but not
benchmarked; write operations face different challenges (write amplification,
garbage collection) that may affect latency profiles. (3)~All measurements
use sequential access patterns; random I/O workloads relevant to graph
analytics~\cite{qureshi2023bam,park2024gids} remain unevaluated. (4)~The
NVIDIA driver patch (\S\ref{sec:impl}) ties the system to specific
kernel and driver versions, limiting deployment ease.

% ============================================================
% 9. CONCLUSION
% ============================================================
\section{Conclusion}
\label{sec:conclusion}

We presented \sysname{}, the first system (to our knowledge) that demonstrates
GPU-initiated NVMe I/O on consumer hardware. Using a GeForce RTX 3090 on an
AMD B450 platform, a single CUDA thread autonomously manages NVMe I/O
queues---constructing commands, writing doorbells via PTX MMIO instructions, and
polling completions---without any CPU involvement in the data path.

Our key insight is that PCIe posted writes (the only operation needed for NVMe
doorbell ringing) succeed on AMD consumer platforms even where P2P reads fail.
This enables a Tier~1 approach where the GPU drives the NVMe controller while
all data flows through host pinned memory.

\sysname{} achieves 2{,}666\MBs{} on the SN530 (78\% of link bandwidth) and
3{,}350\MBs{} on the SN740 (99\% of link bandwidth), outperforming
CPU-mediated methods by up to 2.2$\times$ despite the CPU baselines using a
faster NVMe device. Integrated into a streaming inference engine, \sysname{}
delivers 0.06~tok/s for Llama~70B from NVMe alone, and 0.50~tok/s when
combined with tiered caching and quantization optimizations.

The system demonstrates that GPU-initiated storage I/O is not limited to
enterprise hardware---consumer GPUs are physically capable of this operation,
and the main barriers are software (driver checks, P2P topology) rather than
hardware. This work opens the door to GPU-centric storage access on commodity
hardware, with immediate applications in LLM inference, graph analytics, and
any workload where GPU memory is insufficient to hold the full working set.

% ============================================================
% REFERENCES
% ============================================================
\begin{thebibliography}{99}

\bibitem{qureshi2023bam}
Z.~Qureshi, V.~S.~Mailthody, I.~Gelado, S.~W.~Min, A.~Masood, J.~Park,
J.~Xiong, C.~J.~Newburn, D.~Lopatenko, and W.~Hwu.
\newblock {BaM}: A Case for Enabling Fine-grain High Throughput
  {GPU}-Orchestrated Access to Storage.
\newblock In {\em Proc.\ ASPLOS}, 2023.

\bibitem{gds}
{NVIDIA Corporation}.
\newblock {GPUDirect Storage}.
\newblock \url{https://docs.nvidia.com/gpudirect-storage/}, 2024.

\bibitem{spdk}
Z.~Yang, J.~R.~Harris, B.~Walker, D.~Verkamp, C.~Liu, C.~Chang, G.~Cao,
J.~Stern, V.~Verma, and L.~E.~Paul.
\newblock {SPDK}: A Development Kit to Build High Performance Storage
  Applications.
\newblock In {\em Proc.\ IEEE CloudCom}, 2017.

\bibitem{markussen2021smartio}
J.~Markussen, L.~B.~Kristiansen, P.~Halvorsen, H.~Kielland-Gyrud,
H.~K.~Stensland, and C.~Griwodz.
\newblock {SmartIO}: Zero-overhead Device Sharing through {PCIe} Networking.
\newblock {\em ACM Trans.\ Comput.\ Syst.}, 2021.

\bibitem{silberstein2014gpufs}
M.~Silberstein, B.~Ford, I.~Keidar, and E.~Witchel.
\newblock {GPUfs}: Integrating a File System with {GPUs}.
\newblock {\em ACM Trans.\ Comput.\ Syst.}, 32(1), 2014.

\bibitem{bergman2017spin}
S.~Bergman, T.~Brokhman, T.~Cohen, and M.~Silberstein.
\newblock {SPIN}: Seamless Operating System Integration of Peer-to-Peer {DMA}
  Between {SSDs} and {GPUs}.
\newblock In {\em Proc.\ USENIX ATC}, 2017.

\bibitem{markthub2018dragon}
P.~Markthub, M.~E.~Belviranli, S.~Lee, J.~S.~Vetter, and S.~Matsuoka.
\newblock {DRAGON}: Breaking {GPU} Memory Capacity Limits with Direct {NVM}
  Access.
\newblock In {\em Proc.\ SC}, 2018.

\bibitem{sheng2023flexgen}
Y.~Sheng, L.~Zheng, B.~Yuan, Z.~Li, M.~Ryabinin, B.~Chen, P.~Liang,
C.~Re, I.~Stoica, and C.~Zhang.
\newblock {FlexGen}: High-Throughput Generative Inference of Large Language
  Models with a Single {GPU}.
\newblock In {\em Proc.\ ICML}, 2023.

\bibitem{aminabadi2022deepspeed}
R.~Y.~Aminabadi, S.~Rajbhandari, A.~A.~Awan, C.~Li, D.~Li, E.~Zheng,
O.~Ruwase, S.~Smith, M.~Zhang, J.~Rasley, and Y.~He.
\newblock {DeepSpeed-Inference}: Enabling Efficient Inference of Transformer
  Models at Unprecedented Scale.
\newblock In {\em Proc.\ SC}, 2022.

\bibitem{kwon2023vllm}
W.~Kwon, Z.~Li, S.~Zhuang, Y.~Sheng, L.~Zheng, C.~H.~Yu,
J.~E.~Gonzalez, H.~Zhang, and I.~Stoica.
\newblock Efficient Memory Management for Large Language Model Serving with
  {PagedAttention}.
\newblock In {\em Proc.\ SOSP}, 2023.

\bibitem{song2023powerinfer}
Y.~Song, Z.~Mi, H.~Xie, and H.~Chen.
\newblock {PowerInfer}: Fast Large Language Model Serving with a Consumer-grade
  {GPU}.
\newblock In {\em Proc.\ SOSP}, 2024.

\bibitem{alizadeh2023llmflash}
K.~Alizadeh, I.~Mirzadeh, D.~Belenko, K.~Khatamifard, M.~Cho,
C.~C.~Del~Mundo, M.~Rastegari, and M.~Farajtabar.
\newblock {LLM} in a Flash: Efficient Large Language Model Inference with
  Limited Memory.
\newblock {\em arXiv:2312.11514}, 2023.

\bibitem{park2024gids}
J.~B.~Park, V.~S.~Mailthody, Z.~Qureshi, and W.~Hwu.
\newblock Accelerating Sampling and Aggregation Operations in {GNN} Frameworks
  with {GPU} Initiated Direct Storage Accesses.
\newblock {\em Proc.\ VLDB Endow.}, 17(6):1227--1240, 2024.

\bibitem{fu2024serverlessllm}
Y.~Fu, L.~Xue, Y.~Huang, A.-O.~Brabete, D.~Ustiugov, Y.~Patel, and L.~Mai.
\newblock {ServerlessLLM}: Low-Latency Serverless Inference for Large Language
  Models.
\newblock In {\em Proc.\ OSDI}, 2024.

\bibitem{liu2025pipo}
Y.~Liu, J.~Li, and W.-J.~Li.
\newblock {PIPO}: Pipelined Offloading for Efficient Inference on Consumer
  Devices.
\newblock {\em arXiv:2504.03664}, 2025.

\bibitem{luo2025headinfer}
C.~Luo, C.~Cai, H.~Sun, J.~Xiao, B.~Yuan, W.~Xiao, J.~Hu, J.~Zhao,
B.~Chen, and A.~Anandkumar.
\newblock {HeadInfer}: Memory-Efficient {LLM} Inference by Head-wise
  Offloading.
\newblock {\em arXiv:2502.12574}, 2025.

\bibitem{instinfer2024}
{InstInfer}: In-Storage Attention Offloading for Cost-Effective Long-Context
  {LLM} Inference.
\newblock {\em arXiv:2409.04992}, 2024.

\bibitem{chong2026nanoquant}
H.~Chong, D.~Kim, C.~Kim, and M.~Choi.
\newblock {NanoQuant}: Efficient Sub-1-Bit Quantization of Large Language
  Models.
\newblock {\em arXiv:2602.06694}, 2026.

\bibitem{cheops2025io}
{An I/O Characterizing Study of Offloading LLM Models and KV Caches to NVMe
  SSD}.
\newblock In {\em Proc.\ CHEOPS}, ACM, 2025.

\bibitem{phoenix2025}
{Phoenix}: A Refactored {I/O} Stack for {GPU} Direct Storage without Phony
  Buffers.
\newblock In {\em Proc.\ EuroSys}, ACM, 2025.

\bibitem{zhong2025cxl}
Y.~Zhong \emph{et al.}
\newblock My {CXL} Pool Obviates Your {PCIe} Switch.
\newblock In {\em Proc.\ OSDI}, 2025.

\bibitem{nvme-spec}
{NVM Express Workgroup}.
\newblock NVM Express Base Specification, Revision 2.0.
\newblock \url{https://nvmexpress.org/specifications/}, 2021.

\bibitem{touvron2023llama2}
H.~Touvron, L.~Martin, K.~Stone, P.~Albert, \emph{et al.}
\newblock {Llama 2}: Open Foundation and Fine-Tuned Chat Models.
\newblock {\em arXiv:2307.09288}, 2023.

\end{thebibliography}

\end{document}
